\section{EESQLBench}
\label{sec:EESQLBench}

In this section, we introduce the construction of \benchmark , the first benchmark for evaluating Efficiency-Oriented Text-to-SQL tasks.
Figure \ref{fig:overview} illustrates the entire process involved in building our benchmark.
Starting from real-world databases, we enhance database schemas and data instances to expose performance-sensitive behaviors, and generate   complex SQL queries that admit substantial optimization opportunities.
For each generated query, we further curate a semantically equivalent but efficiency-optimized canonical SQL via automated rewriting and expert refinement, and align it with a natural language (NL) question.
More details will be introduced in the following subsections.

\subsection{Benchmark Construction}
% Implementation

\begin{figure}[t]
	\centering
	\includegraphics[width=0.9\linewidth]{./images/Overview.pdf} 
	\vspace{-5mm}
	\caption{\textbf{Overview of \benchmark\ Construction Pipeline}}
	\label{fig:overview}
	\vspace{-5mm}
\end{figure}

\subsubsection{Database Preparation.}
We prepare the databases of \benchmark\ by building upon real-world relational databases collected from the \textsc{BIRD}~\cite{li2023can} benchmark.
We choose \textsc{BIRD}~\cite{li2023can} because it covers cross-domain scenarios spanning diverse application areas, including e-commerce, finance, and healthcare, and provides large-scale tables derived from real-world data with rich relational structures and realistic data distributions.
Based on these databases, we enhance the original schemas to better support efficiency-oriented evaluation.
% 修复外键
In particular, we explicitly repair missing or implicit foreign-key relationships, which enables multiple alternative join paths among related tables and allows semantically equivalent SQL queries to be expressed in structurally different forms.
% 随机添加索引
We further apply mixed indexing strategies by randomly indexing a subset of foreign-key and attribute columns, so that different join orders and access paths can lead to distinct execution plans. 
% 插入数值和噪声数据，数据分布
At the instance level, we enlarge the databases by systematically inserting additional rows into existing tables while strictly preserving the original schemas, data types, and foreign-key constraints.
For each table, new tuples are generated by sampling valid key values from referenced tables and duplicating existing records with controlled variations on non-key attributes.
We explicitly control value frequencies for selected predicate-relevant attributes by making some values significantly more frequent than others, thereby introducing skewed data distributions that amplify efficiency differences among semantically equivalent queries.


\subsubsection{Complexity-Aware Query Generation.}
After preparing databases, we synthesize SQL queries that are both structurally complex and optimization-sensitive.
We adopt a reverse construction strategy: instead of starting from NL and hoping a generator produces sufficiently challenging SQL, we first generate complex SQL candidates and then align them with NL questions.
%Our design is inspired by DBMS fuzzing, which is effective at exploring a wide spectrum of query structures (e.g., deep nesting, multi-join chains, and aggregation patterns) that can expose optimizer behaviors.
%However, directly using off-the-shelf fuzzers often yields many ill-typed, semantically invalid, or practically meaningless queries, while directly prompting an LLM tends to produce overly ``clean'' and repetitive SQL due to pretraining biases.
%This means LLM-generated queries often exhibit limited structural diversity, favoring simple join patterns and shallow nesting that provide few opportunities for efficiency-sensitive analysis.
Our design is inspired by DBMS fuzzing, which can explore a wide spectrum of query structures (e.g., deep nesting, multi-join chains, and aggregation patterns) that are useful for stress-testing optimizer behavior.
However, off-the-shelf fuzzers are not directly applicable: they often produce ill-typed, semantically invalid, or practically meaningless queries.
On the other hand, prompting an LLM can improve semantic plausibility and executability, but the generated SQL is often overly ``clean'' and repetitive due to pretraining biases, favoring simple join patterns and shallow nesting.
As a result, neither approach alone can simultaneously provide high structural diversity and semantic validity.
To balance these two objectives, we introduce a \emph{SQL Sketch} mechanism that decouples \emph{query structure} from \emph{schema-specific semantics}.
As illustrated in Figure~\ref{fig:sql template}, query generation proceeds in two stages: (i) grammar-based SQL sketch generation, which generates the high-level relational structure of a query, and (ii) LLM-based instantiation, which binds schema-dependent elements under explicit constraints.

\begin{figure}[t]
	\centering
	\includegraphics[width=1\linewidth]{./images/SQL template.pdf} 
	\vspace{-8mm}
	\caption{\textbf{Two-Stage Query Generation via Grammar-Based SQL Sketching and LLM-Based Instantiation}}
	\label{fig:sql template}
	\vspace{-7mm}
\end{figure}

\textbf{(1) Grammar-Based SQL Sketch Generation.}
We firstly generate SQL sketches by sampling from a context-free grammar that specifies optimization-sensitive query structures.
We implement this grammar as shown in Figure \ref{fig:sql template} and use \textsc{go-randgen}~\cite{go-randgen} as a grammar-based generator to derive grammar-valid sketches, covering patterns such as multi-table joins, nested subqueries, common table expressions (CTEs), and aggregations with \texttt{HAVING} clauses.
By construction, each sketch fixes the high-level composition and nesting of relational operators, while abstracting away schema-specific details.
Concretely, schema-dependent elements—including table names, column references, join predicates, and literal constants—are replaced with typed placeholders (denoted by identifiers prefixed with an underscore, e.g., \texttt{\_table}, \texttt{\_cte\_table}, \texttt{\_column}, \texttt{\_agg\_func}, \texttt{\_const}).
This sketch-based design ensures structural correctness and diversity at the query level, while deferring schema binding and semantic instantiation to a later stage.

\textbf{(2) LLM-based Query Instantiation.}
Given a target database, we prompt an LLM to instantiate each sketch by filling placeholders with concrete tables, columns, predicates, and values.
To improve semantic plausibility and controllability, the instantiation prompt includes the following components:
\begin{itemize}[leftmargin=10pt]
	\item \textbf{Task Instruction:} Instructs the LLM to fill the provided SQL sketch into an executable query that corresponds to a meaningful analysis intent, rather than producing arbitrary or trivial SQL.
	\item \textbf{Database Schema:} Provides the full \texttt{CREATE TABLE} statements (table names, columns, types, primary keys, and foreign keys) so that the LLM can select type-consistent columns and produce join conditions aligned with the schema constraints.
	\item \textbf{Database Values:} Samples representative values for a few columns (e.g., categorical values, typical timestamps, common IDs) to ground predicate generation, reducing ill-formed conditions (e.g., comparing strings to numeric columns) and improving contextual relevance.
%	\item \textbf{Column Selection Constraint:} Specifies the target number of output columns (e.g., sampled from a small-biased distribution) so that the generated query returns a concise result schema, consistent with common user-facing analytical questions that focus on a few key attributes.
	\item \textbf{Few-shot Examples:} Provides a small number of example queries following the same formatting and style, helping the LLM adhere to the desired SQL conventions (e.g., aliasing, indentation, and aggregation style).
\end{itemize}
The LLM is explicitly instructed to satisfy typing constraints and foreign-key connectivity when instantiating placeholders, producing executable SQL candidates.

Finally, we validate each instantiated query through execution.
Queries that fail to parse, violate schema constraints, or produce empty results are discarded.
%Only queries that execute successfully and exhibit non-trivial intermediate data processing are retained.
This process results in a large pool of complex and semantically meaningful SQL queries that admit substantial optimization opportunities, which are used in subsequent canonicalization and question alignment.

\subsubsection{Efficient Query Construction.}
For each synthesized SQL query, our goal is to construct a semantically equivalent but execution-efficient canonical query that represents expert-level optimization.
To this end, we adopt a two-stage optimization pipeline that combines automated query rewriting with expert refinement, and explicitly annotate the applied optimization types.
As a first step, we apply an automated query rewriting procedure to identify candidate efficiency improvements.
Specifically, we leverage \textsc{WeTune}~\cite{wang2022wetune}, a rule-based query rewriting framework, to generate semantically equivalent query variants through systematic application of rewriting rules.
For each original query, we evaluate the rewritten variants and retain those that exhibit reduced execution cost.
Building upon the automatically rewritten candidates, database experts further refine the queries to produce a final optimized canonical SQL.
% 专家的选择
In this project, we engage three full-time Text-to-SQL experts, all of whom are Ph.D. students with formal training in hands-on experience in SQL optimization and query tuning. 
The entire expert refinement process for all queries takes approximately 80 person-hours in total.
During this stage, experts manually inspect query structures and execution plans, and apply domain knowledge to introduce deeper optimizations that are difficult to fully capture by rule-based systems.
In addition, we explicitly annotate the primary optimization types applied to each query pair.
Each query pair is labeled with one or more optimization categories indicating where efficiency gains are achieved, such as \emph{join reordering}, \emph{predicate pushdown}, \emph{subquery elimination}, and \emph{aggregation rewriting}.
For example, if an optimized query achieves lower execution cost by pushing a filtering predicate from an outer query into an inner subquery before aggregation, it is annotated with \emph{predicate pushdown}.
These annotations enable fine-grained analysis of Text-to-SQL systems’ optimization capabilities.


\subsubsection{Natural Language Question Alignment.}
After constructing pairs of semantically equivalent original and optimized SQL queries, we align them with NL questions that faithfully capture the underlying analytical intent.
The key design principle of this stage is to ensure fairness: each NL question should reflect the high-level user intent of the query, without encoding any SQL-specific structures or optimization decisions.
To achieve this, we generate NL questions based on the shared semantics of each original–optimized query pair, rather than directly mirroring either concrete SQL formulation.
Specifically, we adopt a back-translation strategy in which an LLM is prompted to translate a given SQL query into a concise and natural NL question.
Importantly, the prompt provides both the original SQL and the optimized canonical SQL as semantic references, and explicitly instructs the model to generate a question that is simultaneously satisfied by both queries.
Concretely, the prompt is constructed with the following components:
\begin{itemize}[leftmargin=10pt]
	\item \textbf{Task Instruction.} 
	The model is instructed to generate a NL question that describes the analytical goal of the queries—i.e., what information is requested—rather than how the result is computed.
	
	\item \textbf{Dual Semantic References.} 
	Both the original SQL query and its optimized canonical counterpart are provided as semantic references.
	The model is instructed to infer the common analytical intent underlying the two queries and generate a single NL question that is valid for both, treating the SQLs as black-box specifications of the result.
	
	\item \textbf{Implementation Abstraction Constraint.} 
	The prompt explicitly prohibits mentioning SQL-specific constructs, such as joins, subqueries, aggregation operators, execution order, or optimization strategies, and requires the question to be phrased in natural, user-centric terms.
	
	\item \textbf{Semantic Equivalence Constraint.} 
	The generated question must remain answerable by any semantically equivalent SQL formulation, ensuring neutrality with respect to different query structures and execution strategies.
\end{itemize}
Following generation, all NL questions undergo a human verification and refinement step.
Three experts check each question for semantic alignment with both SQL variants, clarity, and naturalness, and revise ambiguous or overly verbose descriptions when necessary.
As a result, each NL question in \benchmark\ represents a realistic, user-facing query intent, and that can be correctly answered by multiple semantically equivalent SQL queries.
This alignment strategy enables \benchmark\ to simultaneously evaluate (i) semantic correctness, by checking whether a model-generated SQL answers the NL question correctly, and (ii) efficiency awareness, by comparing the execution efficiency of the generated SQL against the expert-optimized canonical query.

\subsubsection{Quality Control and Annotations.}
\label{subsubsec:quality control}
In the final stage, we conduct quality control for each Text-to-SQL task to ensure the correctness, consistency, and reliability of the benchmark.
This stage consists of semantic consistency verification and task difficulty annotation.

\textbf{(1) Semantic consistency verification.}
We first verify the semantic consistency of each task across both SQL queries and the corresponding NL question.
Specifically, we check that the original SQL query and the optimized SQL query are semantically equivalent by executing them under the same database configuration and comparing their result sets.
Meanwhile, we require the optimized query to consistently achieve lower execution cost, ensuring the presence of a genuine and stable efficiency gap.
We further verify that the NL question can be correctly answered by both SQL queries, confirming that it accurately captures their shared analytical intent.

\textbf{(2) Task difficulty annotation.}
After passing semantic consistency verification, we annotate the difficulty of each Text-to-SQL instance using a rule-based human scoring scheme along three complementary dimensions:
\begin{itemize}[leftmargin=10pt]
	\item \textbf{Schema Complexity:} the number of involved tables, join relationships, and foreign-key dependencies;
	\item \textbf{Question Understanding:} the presence of multiple conditions, implicit constraints, or required reasoning in the NL question;
	\item \textbf{SQL Complexity:} the structural and operator-level complexity of the query, such as subqueries, aggregations, \texttt{HAVING} clauses, and CTEs.
	%	SQL 复杂性是根据 Optimized 的判断的
\end{itemize}
Each dimension is graded according to predefined criteria specified in the annotation guidelines~\cite{eesqlbench_guidelines}, and the final difficulty label (e.g., Simple, Medium, or Hard) is determined by aggregating the scores across all dimensions.

All verification and annotation procedures are conducted by three domain experts.
Two experts independently evaluate each instance following a unified guideline, and disagreements are resolved by a third expert acting as an adjudicator.
We measure inter-annotator agreement using Cohen’s kappa coefficient, obtaining an overall kappa score of 0.92, which indicates high annotation consistency and reliable quality.
As a result, each task in \benchmark\ is accompanied by structured annotations, including the database context, paired original and optimized SQL queries, a shared NL question, explicit optimization-type labels, and a difficulty level.
These annotations provide a reliable foundation for analyzing both semantic correctness and optimization awareness of Text-to-SQL systems.

\subsection{Statistics}

\begin{wrapfigure}{r}{0.45\textwidth}
	\vspace{-14mm}
	\centering
	\includegraphics[width=0.45\textwidth]{./images/DB Complexity.pdf}
	\vspace{-5mm}
	\caption{\textbf{Distributions of Database Size in \benchmark\ (Log Scale)}}
	\label{fig:statistics}
	\vspace{-4mm}
\end{wrapfigure}
We present a statistical analysis of \benchmark\ and compare it with representative Text-to-SQL benchmarks in Table~\ref{tab:text2sql_stats}.
\benchmark\ is constructed on MySQL as the execution backend.
Overall, \benchmark\ contains \tasknumber\ Text-to-SQL tasks over 69 real-world databases.
All tasks are manually curated and paired with expert-optimized canonical SQL queries, enabling direct efficiency comparison between model-generated SQL and expert-level solutions.
Following a unified guideline (Section~\ref{subsubsec:quality control}), we further label each task with a difficulty level, resulting in 55 Simple, 85 Medium, and 73 Hard tasks.

%Figure~\ref{fig:statistics} shows the distribution of database sizes in \benchmark\ (log scale).
Figure~\ref{fig:statistics} shows the distribution of database sizes in \benchmark\ on a log scale; for completeness, we provide a detailed breakdown of record counts in our online documentation~\cite{eesqlbench_database_size}.
\benchmark\ covers databases with highly heterogeneous instance scales, spanning from 148kB to 4.4 GB.
In addition, \benchmark\ is explicitly designed to be efficiency-aware at three levels.
At the \emph{schema level}, each database contains 7.75 tables on average (and up to 65), with repaired foreign-key relationships that enable multiple semantically equivalent join paths.
At the \emph{instance level}, each database contains 3,432,599 records on average, and we further introduce controlled table-scale imbalance and value-frequency skew to amplify cost gaps among semantically equivalent queries.
At the \emph{query level}, each task involves 3.55 tables on average (up to 11) and 2.88 subqueries on average (up to 8), resulting in a substantial portion of multi-join and nested queries that admit non-trivial optimization opportunities.



\section{Efficiency-Oriented Evaluation Metrics}
\label{sec:metrics}
%1、EX
%2、在EX通过的基础上统计 Execution Time 、Actual Rows？
%3、计算与人工解的比值？

Evaluating Efficiency-Oriented Text-to-SQL requires metrics that not only assess semantic correctness but also reliably quantify query execution efficiency. 
While execution latency is a natural indicator of efficiency, it is highly sensitive to external factors, making it unstable and difficult to reproduce. 
Therefore, we adopt query-plan-based cost metrics, which provide a more stable proxy for intrinsic query efficiency.

\paragraph{Semantic Correctness.}
We first evaluate whether a generated query $\hat{q}$ is semantically equivalent to the canonical query $q^*$.
Following prior works~\cite{yu2018spider,li2023can}, semantic equivalence is determined by comparing execution results:
\[
\mathbf{1}(q^*, \hat{q}) =
\begin{cases}
	1, & \text{if } \text{Exec}(\hat{q}, \mathcal{S}) = \text{Exec}(q^*, \mathcal{S}), \\
	0, & \text{otherwise}.
\end{cases}
\]
Efficiency-related metrics are computed only on semantically correct queries (i.e., those with $\mathbf{1}(q^*, \hat{q}) = 1$), ensuring that efficiency evaluation is not confounded by incorrect query semantics.

\paragraph{\textbf{Cost Reachability (CR)}}
To measure a model-generated query's proximity to expert-level efficiency,
we introduce \emph{Cost Reachability (CR)}.
For each semantically correct query pair $(q^*, \hat{q})$, we define:
\[
CR_n = \frac{\text{Cost}(q^*, \mathcal{S})}{\text{Cost}(\hat{q}, \mathcal{S})},
\]
where $\text{Cost}(q, \mathcal{S})$ denotes the execution cost of query $q$ on schema $\mathcal{S}$, 
obtained from the query execution plan.
A larger $CR_n$ indicates that the generated query $\hat{q}$ is more efficient relative to the canonical solution $q^*$, 
while $CR_n < 1$ implies inferior efficiency.

To mitigate the influence of extreme outliers, we follow prior cost-based evaluation practices~\cite{li2023can} and apply a square-root transformation to obtain a stabilized score:
\[
R_n = \sqrt{CR_n}.
\]
The overall Cost Reachability is then computed as:
\[
CR_{\text{avg}} =
\frac{1}{\sum_{n=1}^{N} \mathbf{1}(q^*_n, \hat{q}_n)}
\sum_{n=1}^{N} \mathbf{1}(q^*_n, \hat{q}_n) \cdot R_n,
\]
where $N$ is the total number of evaluation instances. For brevity, we use CR to refer to the aggregated score $CR_{\text{avg}}$ in the remainder of the paper.

\paragraph{\textbf{Acceptable Reachability at $k$ (AR@$k$).}}
While $CR_{\text{avg}}$ characterizes the efficiency gap \emph{conditioned on semantic correctness}, it does not reflect how often a model produces queries that are both correct and practically efficient.
To capture this end-to-end utility, we define \emph{Acceptable Reachability at $k$} (AR@$k$) as:
\[
AR@k =
\frac{1}{N}
\sum_{n=1}^{N} \mathbf{1}(q^*_n, \hat{q}_n) \cdot \mathbf{1}(CR_n \ge k),
\]
where $k \in (0,1]$ is a user-defined threshold.
Intuitively, AR@$k$ measures the fraction of all test instances for which the generated SQL is semantically correct and its execution cost does not exceed $1/k$ times that of the canonical solution.
For example, $k=0.5$ corresponds to accepting queries that are at most twice as expensive as the expert-optimized SQL, while $k=1.0$ only counts queries whose execution cost is no higher than the canonical solution (i.e., at least as efficient as the expert-optimized SQL).

Together, CR and AR@$k$ provide complementary perspectives on efficiency:
CR measures \emph{how close} a model can get to expert-level efficiency on average (given correctness),
while AR@$k$ measures \emph{how often} it produces SQL that is both correct and meets a practical efficiency threshold.

